\begin{proof}[Proof of \cref{obj:lem:dense-moment-matching-lower}]
Let
\[
L_d=\log(ed),\qquad
K_d^{\mathrm{lb}}=2\lceil 4L_d\rceil,\qquad
\lambda_{n,d}=\frac{2n}{d},\qquad
a_{n,d}=\sqrt{\frac{K_d^{\mathrm{lb}}}{64\lambda_{n,d}}}.
\]
By \cref{obj:lem:dense-amplitude-range}, the standing conditions \(2\le d\) and \(d^2<n\) give
\[
0<a_{n,d}\le \frac12.
\]
The ceiling bounds also give
\[
8L_d\le K_d^{\mathrm{lb}}<8L_d+2.
\]
If \(d\ge2\), then \(L_d=1+\log d\le d\). Hence
\[
4L_d\le 4d,\qquad \lceil 4L_d\rceil\le 4d,
\]
because \(4d\) is an integer above \(4L_d\). Therefore \(K_d^{\mathrm{lb}}\le8d\). If also \(d^2<n\), then \(\lambda_{n,d}>0\), and
\[
a_{n,d}^2
=
\frac{K_d^{\mathrm{lb}}}{64\lambda_{n,d}}
=
\frac{K_d^{\mathrm{lb}}d}{128n}.
\]
These relations, together with the amplitude range supplied above, prove the first group of claims.

For \(\theta\in[-a_{n,d},a_{n,d}]^d\), the amplitude bound places every coordinate in \([-1/2,1/2]\). The dense law of \cref{obj:def:dense-moment-matching-construction} has, in every cell,
\[
p_x=\frac1d,\qquad
\pi_x=\frac12,\qquad
\mu_{1x}=\frac{1+\theta_x}{2},\qquad
\mu_{0x}=\frac{1-\theta_x}{2},
\]
equivalently
\[
q_{11,x}=q_{00,x}=\frac{1+\theta_x}{4d},
\qquad
q_{10,x}=q_{01,x}=\frac{1-\theta_x}{4d}.
\]
Since \(0<\epsilon<1/2\), the propensity \(1/2\) belongs to \([\epsilon,1-\epsilon]\), and the displayed means belong to \([0,1]\). Thus the observed margin lies in the class of \cref{obj:def:observed-model-class}. By \cref{obj:def:observed-optimal-value},
\[
\Psi\{\operatorname{Obs}(P_\theta^{\mathrm{dense}})\}
=
\frac1d\sum_{x=1}^d
\max\left\{\frac{1+\theta_x}{2},\frac{1-\theta_x}{2}\right\}
=
\frac12+\frac1{2d}\sum_{x=1}^d|\theta_x|.
\]
By \citep[Lemma 1 and Section 3.1]{CaiLow2011L1}, for every positive even \(K\) there are symmetric probability measures \(\nu_{0,K}\) and \(\nu_{1,K}\) on \([-1,1]\), with matching moments through degree \(K\), such that
\[
\int |t|\,d\nu_{1,K}(t)-\int |t|\,d\nu_{0,K}(t)=2E_K,
\qquad
\frac{c_{\mathrm{CL}}}{K}\le E_K\le \frac{C_{\mathrm{CL}}}{K}
\]
for universal constants \(0<c_{\mathrm{CL}}<C_{\mathrm{CL}}<\infty\). Choose these priors at \(K=K_d^{\mathrm{lb}}\), take their \(d\)-fold products, and map each coordinate by \(t\mapsto a_{n,d}t\). The preceding display for \(\Psi\) gives
\[
\mathbb E_{\Pi_{j,n,d}}\Psi\{\operatorname{Obs}(P_\theta^{\mathrm{dense}})\}
=
\frac12+\frac{a_{n,d}}{2d}
\sum_{x=1}^d\int |t|\,d\nu_{j,K_d^{\mathrm{lb}}}(t),
\]
and hence
\[
\Delta_{n,d}^{\mathrm{dense}}
=
a_{n,d}E_{K_d^{\mathrm{lb}}}.
\]
The lower Cai--Low bound gives \(E_{K_d^{\mathrm{lb}}}>0\). The zero polynomial is an admissible degree-\(K_d^{\mathrm{lb}}\) approximant to \(|t|\), and
\[
\sup_{t\in[-1,1]}||t|-0|\le1,
\]
so \(E_{K_d^{\mathrm{lb}}}\le1\). Together with \(0<a_{n,d}\le1/2\), this yields
\[
0<\Delta_{n,d}^{\mathrm{dense}}\le1.
\]
The same reciprocal approximation bound supplies a universal cutoff \(D_{\mathrm{scale}}\) such that, for every \(d\ge D_{\mathrm{scale}}\),
\[
32\le d\,E_{K_d^{\mathrm{lb}}}^{2}.
\]
Indeed, \(K_d^{\mathrm{lb}}\le10L_d\) for \(d\ge2\), and \(L_d^2=o(\sqrt d)\). Thus, after increasing \(D_{\mathrm{scale}}\),
\[
(K_d^{\mathrm{lb}})^2\le 100\sqrt e\,\sqrt d,
\qquad
3200\sqrt e\le c_{\mathrm{CL}}^2\sqrt d,
\]
and therefore
\[
dE_{K_d^{\mathrm{lb}}}^2
\ge
d\frac{c_{\mathrm{CL}}^2}{(K_d^{\mathrm{lb}})^2}
\ge
\frac{c_{\mathrm{CL}}^2\sqrt d}{100\sqrt e}
\ge32.
\]
For the one-cell likelihood identity, let \(R_+\) and \(R_-\) be independent \(\operatorname{Pois}(\lambda_{n,d}/2)\) counts under the zero-contrast baseline and define
\[
L_\theta(R_+,R_-)
=
(1+\theta)^{R_+}(1-\theta)^{R_-}.
\]
The Poisson probability-generating identity
\[
\int c^r\,d\operatorname{Pois}(\rho)(r)=\exp\{\rho(c-1)\}
\]
gives, for all \(\theta,\theta'\in\mathbb R\),
\[
E_0[L_\theta L_{\theta'}]
=
\exp\!\left[
\frac{\lambda_{n,d}}2\{(1+\theta)(1+\theta')-1\}
+
\frac{\lambda_{n,d}}2\{(1-\theta)(1-\theta')-1\}
\right]
=
\exp(\lambda_{n,d}\theta\theta').
\]
Let \(\rho_{j,n,d}\) be the one-cell mixture density relative to the zero-contrast sign-count baseline \(Q_{0,n,d}=\operatorname{Pois}(\lambda_{n,d}/2)\otimes\operatorname{Pois}(\lambda_{n,d}/2)\):
\[
\rho_{j,n,d}(r_+,r_-)
=
\int L_{a_{n,d}t}(r_+,r_-)\,d\nu_{j,K_d^{\mathrm{lb}}}(t),
\qquad j\in\{0,1\}.
\]
Expanding the exponential kernel above and using moment agreement through \(K_d^{\mathrm{lb}}\) gives
\[
\int (\rho_{1,n,d}-\rho_{0,n,d})^2\,dQ_{0,n,d}
\le
4\sum_{r=K_d^{\mathrm{lb}}+1}^{\infty}
\frac{(\lambda_{n,d}a_{n,d}^2)^r}{r!}
=
4\,\operatorname{tail}_{n,d}^{\mathrm{dense}}.
\]
Define the supported one-cell sign kernel on \(\mathbb R\) by
\[
\mathsf K_{n,d}(t,\cdot)
=
\begin{cases}
\operatorname{Law}\!\left(
\operatorname{Pois}\{\lambda_{n,d}(1+a_{n,d}t)/2\},
\operatorname{Pois}\{\lambda_{n,d}(1-a_{n,d}t)/2\}
\right), & |t|\le1,\\[0.4ex]
Q_{0,n,d}, & |t|>1,
\end{cases}
\]
with independent coordinates in the first line. Its \(d\)-cell product prior-predictive law is
\[
\mathfrak S_{j,n,d}
=
\int
\bigotimes_{x=1}^d \mathsf K_{n,d}(t_x,\cdot)\,
d\nu_{j,K_d^{\mathrm{lb}}}^{\otimes d}(t_1,\ldots,t_d),
\qquad j\in\{0,1\}.
\]
The Cauchy--Schwarz total-variation bound for one cell and the telescoping inequality for products yield
\[
\operatorname{TV}(\mathfrak S_{0,n,d},\mathfrak S_{1,n,d})
\le
d\{\operatorname{tail}_{n,d}^{\mathrm{dense}}\}^{1/2}.
\]
Let \(\operatorname{Sign}\) map a Poisson observed sample to the array whose \(x\)-th entry is the pair of counts with \(A=Y\) and \(A\ne Y\) in cell \(x\). Under the zero-contrast dense law, let
\[
\mathsf R_{n,d}(r,\cdot)
=
\Pr_{0}\{\text{Poisson observed sample}\in\cdot\mid \operatorname{Sign}=r\}
\]
be a regular conditional law. For every supported product prior above, this reconstruction kernel satisfies
\[
\mathbb M_{j,n,d}(\mathcal B)
=
\int \mathsf R_{n,d}(r,\mathcal B)\,d\mathfrak S_{j,n,d}(r),
\qquad j\in\{0,1\},
\]
for every measurable observed-sample event \(\mathcal B\). Since total variation weakly decreases under a common Markov kernel,
\[
\operatorname{TV}(\mathbb M_{0,n,d},\mathbb M_{1,n,d})
\le
\operatorname{TV}(\mathfrak S_{0,n,d},\mathfrak S_{1,n,d})
\le
d\{\operatorname{tail}_{n,d}^{\mathrm{dense}}\}^{1/2}.
\]
Since
\[
\lambda_{n,d}a_{n,d}^2=\frac{K_d^{\mathrm{lb}}}{64},
\]
Stirling's lower bound implies, for every \(r\ge K_d^{\mathrm{lb}}+1\),
\[
\frac{(\lambda_{n,d}a_{n,d}^2)^r}{r!}
=
\frac{(K_d^{\mathrm{lb}}/64)^r}{r!}
\le
\left(\frac e{64}\right)^r.
\]
Therefore
\[
\operatorname{tail}_{n,d}^{\mathrm{dense}}
\le
\sum_{r=K_d^{\mathrm{lb}}+1}^\infty \left(\frac e{64}\right)^r
=
\frac{(e/64)^{K_d^{\mathrm{lb}}+1}}{1-e/64}.
\]
Put \(\zeta_{\mathrm{geom}}=e/64\). Since \(\zeta_{\mathrm{geom}}\le e^{-2}\), \(K_d^{\mathrm{lb}}\ge8L_d\), and \(L_d=\log(ed)\),
\[
\zeta_{\mathrm{geom}}^{K_d^{\mathrm{lb}}+1}
\le
e^{-2(K_d^{\mathrm{lb}}+1)}
\le
e^{-16L_d}
=
(ed)^{-16}
\le
\frac{1}{512d^2}
\qquad(d\ge2).
\]
Also \(1-\zeta_{\mathrm{geom}}\ge1/2\), and so
\[
\frac{\zeta_{\mathrm{geom}}^{K_d^{\mathrm{lb}}+1}}{1-\zeta_{\mathrm{geom}}}
\le
\frac{1}{256d^2},
\qquad
d\left\{\frac{\zeta_{\mathrm{geom}}^{K_d^{\mathrm{lb}}+1}}{1-\zeta_{\mathrm{geom}}}\right\}^{1/2}
\le
\frac1{16}.
\]
Combining the last displays proves the total-variation chain in the statement.

Under either product prior, write \(\theta_x=a_{n,d}\Xi_x\), where \(\Xi_1,\ldots,\Xi_d\) are independent with common law \(\nu_{j,K_d^{\mathrm{lb}}}\). The support condition makes the restriction to \([-1,1]^d\) immaterial, and
\[
\Psi\{\operatorname{Obs}(P_\theta^{\mathrm{dense}})\}
=
\frac12+\frac{a_{n,d}}{2d}\sum_{x=1}^d|\Xi_x|.
\]
Because \(|\Xi_x|\in[0,1]\), its variance is at most \(1/4\). Independence gives
\[
\int
\left(\Psi\{\operatorname{Obs}(P_\theta^{\mathrm{dense}})\}
-\mathbb E_{\Pi_{j,n,d}}\Psi\{\operatorname{Obs}(P_\theta^{\mathrm{dense}})\}\right)^2
\,d\Pi_{j,n,d}(\theta)
\le
\frac{a_{n,d}^2}{4d}.
\]
Chebyshev's inequality, \(\Delta_{n,d}^{\mathrm{dense}}=a_{n,d}E_{K_d^{\mathrm{lb}}}\), and \(32\le dE_{K_d^{\mathrm{lb}}}^2\) then yield
\[
\Pi_{j,n,d}\left(
\left|\Psi\{\operatorname{Obs}(P_\theta^{\mathrm{dense}})\}
-\mathbb E_{\Pi_{j,n,d}}\Psi\{\operatorname{Obs}(P_\theta^{\mathrm{dense}})\}\right|
>
\frac{\Delta_{n,d}^{\mathrm{dense}}}{4}
\right)
\le
\frac{a_{n,d}^2/(4d)}
{(a_{n,d}E_{K_d^{\mathrm{lb}}}/4)^2}
=
\frac4{dE_{K_d^{\mathrm{lb}}}^2}
\le
\frac18.
\]
For any estimator in the Poissonized experiment, set
\[
\bar\psi_j
=
\mathbb E_{\Pi_{j,n,d}}
\Psi\{\operatorname{Obs}(P_\theta^{\mathrm{dense}})\},
\qquad
r_\star
=
\max_{j\in\{0,1\}}
\int E_\theta\!\left[
\left\{\widehat V-\Psi\{\operatorname{Obs}(P_\theta^{\mathrm{dense}})\}\right\}^2
\right]d\Pi_{j,n,d}(\theta).
\]
Use the test
\[
\varphi=\mathbf 1\!\left\{\widehat V\ge\frac{\bar\psi_0+\bar\psi_1}{2}\right\},
\qquad
\alpha_0=\mathbb M_{0,n,d}(\varphi=1),
\qquad
\alpha_1=\mathbb M_{1,n,d}(\varphi=0).
\]
On the concentration events above, estimation error smaller than \(\Delta_{n,d}^{\mathrm{dense}}/4\) forces the test to select the correct prior index. Markov's inequality gives
\[
\alpha_0+\alpha_1
\le
\frac14+\frac{32r_\star}{(\Delta_{n,d}^{\mathrm{dense}})^2}.
\]
Every test also satisfies
\[
\alpha_0+\alpha_1
\ge
1-\operatorname{TV}(\mathbb M_{0,n,d},\mathbb M_{1,n,d})
\ge
\frac{15}{16}.
\]
Consequently
\[
\frac{11}{512}(\Delta_{n,d}^{\mathrm{dense}})^2
\le
\mathfrak R^{\mathrm{Pois}}_{2n,d,\epsilon}.
\]
The restriction to dense observed laws is legitimate because those laws lie in \(\mathcal D_{d,\epsilon}^{\mathrm{obs}}\). In addition, \(\Psi(\mathbb P)\in[0,1]\) for every \(\mathbb P\in\mathcal D_{d,\epsilon}^{\mathrm{obs}}\), since \(p_x\ge0\), \(\sum_xp_x=1\), and \(0\le\max\{\mu_{0x},\mu_{1x}\}\le1\); therefore projection of estimators onto \([0,1]\) weakly decreases squared loss.

The Cai--Low lower bound on \(E_K\), together with the amplitude identity, gives a universal constant
\[
c_0=\frac{11c_{\mathrm{CL}}^2}{512\cdot128}>0
\]
such that
\[
c_0\frac{d}{nK_d^{\mathrm{lb}}}
\le
\frac{11}{512}(\Delta_{n,d}^{\mathrm{dense}})^2.
\]
Indeed,
\[
(\Delta_{n,d}^{\mathrm{dense}})^2
=
\frac{K_d^{\mathrm{lb}}d}{128n}E_{K_d^{\mathrm{lb}}}^2
\ge
\frac{c_{\mathrm{CL}}^2}{128}\frac{d}{nK_d^{\mathrm{lb}}}.
\]
Hence
\[
c_0\frac{d}{nK_d^{\mathrm{lb}}}
\le
\mathfrak R^{\mathrm{Pois}}_{2n,d,\epsilon}.
\]
It remains to compare the fixed-sample and Poissonized experiments. Given a fixed-\(n\) estimator, project it onto \([0,1]\), apply it to the first \(n\) observations of a mean-\(2n\) Poisson sample on the event \(\{\operatorname{Pois}(2n)\ge n\}\), and output \(0\) on the complementary event. On the first event, the first \(n\) observations have the \(n\)-fold i.i.d. law from \cref{obj:ass:iid-sampling}; on the second event, the projected squared loss is at most \(1\). Taking infima gives
\[
\mathfrak R_{n,d,\epsilon}
\ge
\mathfrak R^{\mathrm{Pois}}_{2n,d,\epsilon}
-
\Pr\{\operatorname{Pois}(2n)<n\}.
\]
For the lower-tail absorption, fix any \(\eta>0\). Since \(L_d=o(d)\), there is \(D_{\log}\) such that
\[
L_d\le \eta d
\qquad(d\ge D_{\log}).
\]
Increasing the cutoff to ensure \(d\ge20\), the dense condition \(d^2<n\) gives \(n\ge400\). The mean-\(2n\) Poisson Chernoff bound gives
\[
\Pr\{\operatorname{Pois}(2n)<n\}
\le
\exp\{-n(1-\log2)\}.
\]
Since \(1-\log2>3/10\) and \(\log n\le2\sqrt n\) for \(n\ge400\),
\[
2\log n\le n(1-\log2),
\]
and therefore
\[
\exp\{-n(1-\log2)\}\le n^{-2}.
\]
For \(d\ge D_{\log}\) and \(n\ge1\),
\[
n^{-2}\le \frac{\eta d}{nL_d},
\]
because \(L_d\le \eta dn\). Thus, with \(\eta=c_0/20\), there is a cutoff \(D_{\mathrm{tail}}\) such that
\[
\Pr\{\operatorname{Pois}(2n)<n\}
\le
\frac{c_0}{20}\frac{d}{nL_d}
\]
throughout the dense regime above that cutoff.

Choose
\[
D_0=\max\{2,D_{\mathrm{scale}},D_{\mathrm{tail}}\},
\qquad
c_\epsilon=\min\left\{\frac{11}{512},\,c_0,\,\frac{c_0}{20}\right\}.
\]
Then \(D_0\ge2\) and \(c_\epsilon>0\). The construction, observed-submodel identities, moment-matched priors, concentration bounds, total-variation chain, and one-cell likelihood identity have already been established for every \(d\ge D_0\), \(d^2<n\), and \(0<\epsilon<1/2\). Since \(c_\epsilon\le11/512\),
\[
c_\epsilon(\Delta_{n,d}^{\mathrm{dense}})^2
\le
\mathfrak R^{\mathrm{Pois}}_{2n,d,\epsilon}.
\]
Since \(c_\epsilon\le c_0\),
\[
c_\epsilon\frac{d}{nK_d^{\mathrm{lb}}}
\le
\mathfrak R^{\mathrm{Pois}}_{2n,d,\epsilon}.
\]
Finally, \(K_d^{\mathrm{lb}}\le10L_d\), so
\[
\frac{c_0}{10}\frac{d}{nL_d}
\le
c_0\frac{d}{nK_d^{\mathrm{lb}}}
\le
\mathfrak R^{\mathrm{Pois}}_{2n,d,\epsilon}.
\]
Using the Poisson-tail comparison and the fixed-sample transfer,
\[
\frac{c_0}{20}\frac{d}{nL_d}
\le
\mathfrak R^{\mathrm{Pois}}_{2n,d,\epsilon}
-
\Pr\{\operatorname{Pois}(2n)<n\}
\le
\mathfrak R_{n,d,\epsilon}.
\]
Because \(c_\epsilon\le c_0/20\), this proves
\[
c_\epsilon\frac{d}{nL_d}
\le
\mathfrak R_{n,d,\epsilon}.
\]
All displayed conclusions follow.
\end{proof}
